% =====================================================================
% lb_extensions.tex  (Path B: single-segment LB only)
% ---------------------------------------------------------------------
% Drafted upgrade to App.~A.6 (Proof of Theorem~\ref{thm:spsc_lb}).
%
% Addresses the more important of the two reviewer-flagged weaknesses:
%
%   (W2) The current LB is restricted to the hard-projection class
%        $\Pi_{\mathrm{HP}}$, which excludes ambient methods (LinUCB,
%        D-LinUCB, ...) and reads as a self-serving restriction.
%
% Resolution: a tightened KL bound (Lemma~\ref{lem:trace_kl}) using the
% trace identity $\mathbb E_v[v v^\top] = I_d/d$ instead of the
% Cauchy--Schwarz-then-average argument of the original proof. The
% bound becomes policy-agnostic (no $\Pi_{\mathrm{HP}}$ assumption) and
% the constant improves by a factor of $r$.
%
% The companion weakness (W1: single-segment $\to$ multi-segment with
% the joint $r\sqrt{KT}+T^{2/3}$ rate) is acknowledged as an open
% direction. Multiple drafting attempts revealed that a referee-quality
% multi-segment proof requires careful side-observation arguments and
% adaptive-probe-count handling that exceed the scope of this revision.
% This file therefore commits only to the single-segment, policy-class-
% agnostic improvement, and explicitly states the multi-segment
% extension as future work.
%
% INTEGRATION INSTRUCTIONS:
%   * Replace App.~A.6 (lines ~1715--1845 of conf.tex) with the
%     theorem statement and proof in this file.
%   * Update the abstract sentence at lines ~92--96 to reflect the
%     dropped $\Pi_{\mathrm{HP}}$ restriction; do NOT claim $r\sqrt{KT}$
%     matching outside $\Pi_{\mathrm{HP}}$.
%   * Update Theorem~\ref{thm:spsc_lb} statement in main body
%     (lines ~541--558) accordingly; keep the existing $\Pi_{\mathrm{HP}}$
%     LB as a separate "stronger LB inside the projection class"
%     sub-result if desired.
% =====================================================================


\subsection{Sharpened single-segment lower bound (no policy-class restriction)}
\label{app:proof_lb_sharp}

We refine the lower-bound argument of \S\ref{app:proof_lb} via the
trace identity, which removes the hard-projection restriction
$\Pi_{\mathrm{HP}}$ that the original proof imposed and improves the
constant by a factor of $r$.

\begin{lemma}[Trace KL bound, policy-agnostic]
\label{lem:trace_kl}
Fix $T\ge 1$ and $0\le m\le T$, and let $\pi$ be any (possibly
randomized) history-adapted policy taking actions $x_t\in\R^d$ with
$\|x_t\|_2\le R_\cA$ on exploit rounds, and $x_t = u_t$ with
$u_t \sim \mathcal N(0,I_d)$ on the $m$ designated probe rounds (with
$m$ fixed in advance; the bound applies to any deterministic probe
schedule). Let $\cT_{\mathrm{ex}}\subseteq[T]$ denote the exploit
rounds (so $|\cT_{\mathrm{ex}}|=T-m$). Under $\nu_0$, for any random
unit vector $v\sim\mathrm{Unif}(\mathbb S^{d-1})$ chosen
\emph{after} the policy's distribution under $\nu_0$ is fixed,
\begin{equation}
\mathbb E_v\!\left[\,\mathbb E_{\mathbb P_0}\!\Bigl[\sum_{t=1}^T(x_t^\top v)^2\Bigr]\right]
\;\le\; m \;+\; \frac{R_\cA^{\,2}\,(T-m)}{d}.
\label{eq:trace_kl}
\end{equation}
Restricting the inner sum to exploit rounds gives the sharper
\begin{equation}
\mathbb E_v\!\left[\,\mathbb E_{\mathbb P_0}\!\Bigl[\sum_{t\in\cT_{\mathrm{ex}}}(x_t^\top v)^2\Bigr]\right]
\;\le\; \frac{R_\cA^{\,2}\,(T-m)}{d}.
\label{eq:trace_kl_ex}
\end{equation}
Consequently there exists $v^\star\in\mathbb S^{d-1}$ for which the
$\mathbb P_0$-expectations in \eqref{eq:trace_kl} and
\eqref{eq:trace_kl_ex} hold with $v$ replaced by $v^\star$.
\end{lemma}
\begin{proof}
By the trace identity
$\mathbb E_{v\sim\mathrm{Unif}(\mathbb S^{d-1})}[v v^\top] = I_d/d$,
$\mathbb E_v[(x^\top v)^2] = x^\top \mathbb E_v[vv^\top] x = \|x\|^2/d$
for any $x\in\R^d$.

\emph{Probe rounds.} For $u_t\sim\mathcal N(0,I_d)$,
$\mathbb E_{u_t}[\|u_t\|^2]=d$, so
$\mathbb E_v\mathbb E_{u_t}[(u_t^\top v)^2]=\mathbb E_{u_t}[\|u_t\|^2/d]=1$,
contributing $m$ over the $m$ probe rounds.

\emph{Exploit rounds.} On exploit rounds $\|x_t\|\le R_\cA$ a.s., and
the trajectory law under $\mathbb P_0$ is independent of $v$ (rewards
under $\nu_0$ are $\cN(0,\sigma_\varepsilon^2)$ regardless of $v$), so
we may swap the order of expectations by Fubini:
\[
\mathbb E_v\,\mathbb E_{\mathbb P_0}\!\bigl[(x_t^\top v)^2\bigr]
\;=\; \mathbb E_{\mathbb P_0}\!\bigl[\mathbb E_v[(x_t^\top v)^2\mid x_t]\bigr]
\;=\; \mathbb E_{\mathbb P_0}\!\bigl[\|x_t\|^2/d\bigr] \;\le\; R_\cA^2/d,
\]
contributing $R_\cA^2(T-m)/d$ over the $T-m$ exploit rounds.

Adding the two contributions yields \eqref{eq:trace_kl}; restricting
to exploit rounds yields \eqref{eq:trace_kl_ex}.

\emph{A single $v^\star$ achieves both bounds.} Define
$g_{\mathrm{full}}(v) := \mathbb E_{\mathbb P_0}[\sum_{t=1}^T(x_t^\top v)^2]$
and
$g_{\mathrm{ex}}(v) := \mathbb E_{\mathbb P_0}[\sum_{t\in\cT_{\mathrm{ex}}}(x_t^\top v)^2]$.
The difference is the contribution from probe rounds:
$g_{\mathrm{full}}(v) - g_{\mathrm{ex}}(v)
= \mathbb E_{\mathbb P_0}[\sum_{t\in\cT_{\mathrm{probe}}}(u_t^\top v)^2]$.
For every unit $v\in\mathbb S^{d-1}$, $u_t^\top v\sim \cN(0,1)$ (since
$u_t\sim\cN(0,I_d)$ is rotationally invariant), so each probe round
contributes exactly $1$ to this expectation; with $m$ probe rounds,
$g_{\mathrm{full}}(v) - g_{\mathrm{ex}}(v) = m$ for every $v$. Hence
any minimizer $v^\star\in\mathbb S^{d-1}$ of $g_{\mathrm{ex}}$ that
witnesses $g_{\mathrm{ex}}(v^\star)\le \mathbb E_v[g_{\mathrm{ex}}(v)]
\le R_\cA^2(T-m)/d$ \emph{simultaneously} witnesses
$g_{\mathrm{full}}(v^\star) = m + g_{\mathrm{ex}}(v^\star)\le m + R_\cA^2(T-m)/d$,
proving both bounds with the same $v^\star$.
\end{proof}

\begin{remark}[Comparison with the proof of \S\ref{app:proof_lb}]
\label{rem:trace_vs_cs}
The proof of \S\ref{app:proof_lb} bounds exploit-round contributions
via Cauchy--Schwarz:
$(x_t^\top v)^2 = (x_t^\top \widehat P_t v)^2 \le R_\cA^2\|\widehat P_t v\|^2$
(valid only on $\Pi_{\mathrm{HP}}$, where $x_t = \widehat P_t x_t$),
and then averages: $\mathbb E_v\|\widehat P_t v\|^2 = r/d$.
The trace identity sidesteps Cauchy--Schwarz entirely:
$\mathbb E_v[(x_t^\top v)^2] = \|x_t\|^2/d$ holds for arbitrary $x_t$,
so the bound applies to every policy with $\|x_t\|\le R_\cA$.
The factor $r$ between the two bounds is the source of the apparent
$\Pi_{\mathrm{HP}}$ restriction in the earlier proof; with the trace
identity, no such restriction is needed.
\end{remark}

\begin{theorem}[Costed-regret lower bound, all bounded-action policies]
\label{thm:spsc_lb_sharp}
Let $\Pi$ denote the class of all (history-adapted) policies that play
$x_t\in\cA = \{x : \|x\|\le R_\cA\}$ on exploit rounds and isotropic
Gaussian probes $x_t=u_t\sim\mathcal N(0,I_d)$ on probe rounds, with
total probe cost $cm$ included in the dynamic regret. There exist
constants $c_0,d_0,T_0>0$, depending only on
$R_\cA,\sigma_\varepsilon,c$, such that for $d\ge d_0 T^{1/3}$ and
$T\ge T_0$,
\[
\inf_{\pi\in\Pi}\sup_\nu\ \mathbb E^\pi_\nu[\mathrm{DynReg}_T^{(c)}]
\;\ge\; c_0\,c^{1/3}(R_\cA\sigma_\varepsilon)^{2/3}\,T^{2/3}.
\]
\end{theorem}

\begin{proof}
\emph{Construction.} Take $K=1$ (single segment of length $T$). Two
candidate instances:
\begin{itemize}
\item $\nu_0:\theta_t\equiv 0$ for all $t\in[T]$.
\item $\nu_1(v):\theta_t\equiv\alpha v$ for all $t\in[T]$, with $v\in\mathbb S^{d-1}$.
\end{itemize}
Both are valid rank-$1$ piecewise low-rank instances (rank-$r$ with
$r=1$). The scale $\alpha>0$ is set below.

\emph{KL via Lemma~\ref{lem:trace_kl}.} The sequential KL chain rule,
applied across $t=1,\dots,T$, gives
\[
\mathrm{KL}(\mathbb P_0^\cH\|\mathbb P_1^\cH(v))
\;=\; \frac{\alpha^2}{2\sigma_\varepsilon^2}\,
\mathbb E^\pi_{\mathbb P_0}\!\Bigl[\sum_{t=1}^T (x_t^\top v)^2\Bigr].
\]
Picking $v=v^\star$ as in Lemma~\ref{lem:trace_kl} and writing
$Q := m + R_\cA^2(T-m)/d$,
\begin{equation}
\mathrm{KL}(\mathbb P_0^\cH\|\mathbb P_1^\cH(v^\star))
\;\le\; \frac{\alpha^2 Q}{2\sigma_\varepsilon^2}.
\label{eq:singleseg_KL}
\end{equation}

\emph{Le Cam two-point.} Set
$\alpha^2 := \sigma_\varepsilon^2/(4Q)$, so
$\mathrm{KL}\le 1/8$ and Pinsker's inequality gives
$d_{\mathrm{TV}}(\mathbb P_0^\cH,\mathbb P_1^\cH(v^\star))\le 1/4$.

\emph{Cumulative regret on $\nu_1(v^\star)$.} The optimal action under
$\nu_1(v^\star)$ is $x_t^\star = R_\cA v^\star$ with reward
$\alpha R_\cA$. Probe rounds are not comparable to the optimal arm
(their action $u_t$ is fixed by the policy's probe schedule, drawn from
$\cN(0,I_d)$, and unbounded); they enter the dynamic regret only
through the cost $cm$. Per-round arm-regret is therefore
measured only on exploit rounds, where $\|x_t\|\le R_\cA$. The exploit
per-round regret is
$\Delta_t = \alpha\bigl(R_\cA - v^{\star\top} x_t\bigr)
\ge \alpha\bigl(R_\cA - |v^{\star\top} x_t|\bigr)$, and summing,
\begin{equation}
\mathbb E^\pi_{\nu_1(v^\star)}\!\Bigl[\sum_{t\in\cT_{\mathrm{ex}}}\Delta_t\Bigr]
\;\ge\; \alpha (T-m) R_\cA
\;-\; \alpha\,\mathbb E^\pi_{\nu_1(v^\star)}[f],
\label{eq:reg_to_alignment}
\end{equation}
where $f(\cH) := \sum_{t\in\cT_{\mathrm{ex}}} |v^{\star\top}x_t|$.

\emph{TV transfer.} Because $|v^{\star\top}x_t|\le \|x_t\|\le R_\cA$
on every exploit round, $\|f\|_\infty\le (T-m) R_\cA$. (We
\emph{exclude} probe rounds from the sum precisely because
$u_t\sim\cN(0,I_d)$ is unbounded.) For any bounded measurable
functional $f$ of the trajectory $\cH$,
$|\mathbb E^\pi_{\nu_1(v^\star)}[f] - \mathbb E^\pi_{\nu_0}[f]|
\le 2\|f\|_\infty\,d_{\mathrm{TV}}(\mathbb P_0^\cH,\mathbb P_1^\cH(v^\star))$
\citep[Lemma~2.1]{tsybakov2009introduction}, hence with
$d_{\mathrm{TV}}\le 1/4$,
\begin{equation}
\mathbb E^\pi_{\nu_1(v^\star)}[f]
\;\le\; \mathbb E^\pi_{\nu_0}[f]
\;+\; (T-m) R_\cA/2.
\label{eq:tv_transfer}
\end{equation}

\emph{Bound under $\nu_0$.} By Cauchy--Schwarz,
$\mathbb E^\pi_{\nu_0}[f]
\le \sqrt{(T-m)\cdot\mathbb E^\pi_{\nu_0}\!\bigl[\sum_{t\in\cT_{\mathrm{ex}}}(v^{\star\top}x_t)^2\bigr]}$.
Lemma~\ref{lem:trace_kl}, equation \eqref{eq:trace_kl_ex} gives
$\mathbb E^\pi_{\nu_0}[\sum_{t\in\cT_{\mathrm{ex}}}(v^{\star\top}x_t)^2]
\le R_\cA^2(T-m)/d$, so
\begin{equation}
\mathbb E^\pi_{\nu_0}[f]
\;\le\; \sqrt{(T-m)\cdot R_\cA^2(T-m)/d}
\;=\; R_\cA(T-m)/\sqrt d.
\label{eq:nu0_bound}
\end{equation}

\emph{Combining.} Substituting \eqref{eq:tv_transfer} and
\eqref{eq:nu0_bound} into \eqref{eq:reg_to_alignment},
\[
\mathbb E^\pi_{\nu_1(v^\star)}\!\Bigl[\sum_{t\in\cT_{\mathrm{ex}}}\Delta_t\Bigr]
\;\ge\; \alpha (T-m) R_\cA \cdot \Bigl(1 - \frac{1}{\sqrt d} - \frac{1}{2}\Bigr).
\]
For $d\ge 16$, $1/\sqrt d\le 1/4$, so the bracket is at least
$1 - 1/4 - 1/2 = 1/4$:
\begin{equation}
\mathbb E^\pi_{\nu_1(v^\star)}\!\Bigl[\sum_{t\in\cT_{\mathrm{ex}}}\Delta_t\Bigr]
\;\ge\; \tfrac{1}{4}\,\alpha (T-m) R_\cA.
\label{eq:reg_clean}
\end{equation}

\emph{Total dynamic regret.} The dynamic regret is
$\mathrm{DynReg}^{(c)}_T = cm + \sum_{t\in\cT_{\mathrm{ex}}}\Delta_t$,
and substituting $\alpha = \sigma_\varepsilon/(2\sqrt{Q})$ from
\eqref{eq:singleseg_KL},
\begin{equation}
\sup_\nu\mathbb E^\pi_\nu[\mathrm{DynReg}^{(c)}_T]
\;\ge\; cm \;+\; \frac{R_\cA\sigma_\varepsilon\,(T-m)}{8\sqrt{Q}}.
\label{eq:dynreg_lower}
\end{equation}

\emph{Optimization over $m$.} The lower bound in
\eqref{eq:dynreg_lower} must hold for every policy choice of $m$,
which is selected by the policy and may take any value in $\{0,\ldots,T\}$.
We handle the cases $m>T/2$ and $m\le T/2$ separately.

\smallskip\noindent
(\textbf{Preliminary case: $m>T/2$.}) The probe-cost term alone gives
$\mathrm{DynReg}^{(c)}_T \ge cm > cT/2$. For
$T\ge T_0 := 8(R_\cA\sigma_\varepsilon)^{2}/c^{2}$ (a constant
depending only on $R_\cA,\sigma_\varepsilon,c$), the inequality
$cT^{1/3}/2 \ge c^{1/3}(R_\cA\sigma_\varepsilon)^{2/3}$ holds, hence
$cT/2\ge c^{1/3}(R_\cA\sigma_\varepsilon)^{2/3} T^{2/3}$, which already
establishes the claimed rate. Hence in the remainder of the proof we
may assume $m\le T/2$, so that $T-m\ge T/2$.

\smallskip\noindent
(\textbf{Regime I: $R_\cA^2(T-m)/d \le m$, equivalently $Q\le 2m$.})
Then $1/\sqrt{Q}\ge 1/\sqrt{2m}$, and \eqref{eq:dynreg_lower} combined
with $T-m\ge T/2$ gives
$\sup_\nu\mathbb E[\mathrm{DynReg}^{(c)}_T]\ge cm + R_\cA\sigma_\varepsilon T/(16\sqrt{2}\,\sqrt m)$.
Minimizing the RHS over $m$ at
$m^\star\asymp(R_\cA\sigma_\varepsilon T/c)^{2/3}$ yields
\[
\sup_\nu\mathbb E[\mathrm{DynReg}^{(c)}_T]
\;\ge\; c_1\,c^{1/3}(R_\cA\sigma_\varepsilon)^{2/3}T^{2/3}
\]
for an explicit absolute constant $c_1 = 3\cdot 2^{-2/3}/8\approx 0.236$.
The optimum lies inside this regime when
$m^\star\ge R_\cA^2(T-m^\star)/d$, which (for $m^\star\le T/2$) is
implied by $m^\star\ge R_\cA^2 T/d$, i.e., $d\ge R_\cA^2 T/m^\star
\asymp R_\cA^{4/3}(c/\sigma_\varepsilon)^{2/3}\,T^{1/3}$. This is the
condition $d\ge d_0 T^{1/3}$ with $d_0$ depending on
$R_\cA,\sigma_\varepsilon,c$.

\smallskip\noindent
(\textbf{Regime II: $R_\cA^2(T-m)/d > m$.})
Then the second summand of $Q = m + R_\cA^2(T-m)/d$ dominates the
first, and combining with $T-m\le T$ gives the explicit bound
$Q \;<\; 2R_\cA^2(T-m)/d \;\le\; 2R_\cA^2 T/d$.
Therefore $1/\sqrt{Q}\ge \sqrt{d/(2R_\cA^2 T)}$, and
\eqref{eq:dynreg_lower} together with $T-m\ge T/2$ gives
\[
\sup_\nu\mathbb E[\mathrm{DynReg}^{(c)}_T]
\;\ge\; \frac{R_\cA\sigma_\varepsilon (T-m)}{8}\cdot\sqrt{\frac{d}{2R_\cA^2 T}}
\;\ge\; \frac{\sigma_\varepsilon\sqrt{dT}}{16\sqrt 2}.
\]
Under $d\ge d_0 T^{1/3}$, this is at least
$\sigma_\varepsilon\sqrt{d_0}\,T^{2/3}/(16\sqrt 2)$. Choosing $d_0$ large
enough that this lower bound exceeds the Regime~I constant
$c_1\,c^{1/3}(R_\cA\sigma_\varepsilon)^{2/3}$ — concretely,
$\sqrt{d_0}\ge 16\sqrt 2\,c_1\,c^{1/3}R_\cA^{2/3}\sigma_\varepsilon^{-1/3}$,
i.e., $d_0\ge 512\,c_1^2\,c^{2/3}R_\cA^{4/3}\sigma_\varepsilon^{-2/3}$
suffices — Regime~II's bound dominates and the combined LB is
$\Omega(c^{1/3}(R_\cA\sigma_\varepsilon)^{2/3}T^{2/3})$ in both
regimes.

Finally, the bracket simplification at \eqref{eq:reg_clean} used
$d\ge 16$. To absorb this into the global hypothesis $d\ge d_0 T^{1/3}$,
we set the final $d_0$ to the maximum of the Regime-II-domination
constant $512\,c_1^2\,c^{2/3}R_\cA^{4/3}\sigma_\varepsilon^{-2/3}$ and
$16$. Then $d\ge d_0 T^{1/3}\ge d_0\ge 16$ automatically.\qed
\end{proof}

\begin{remark}[Resolution of the $\Pi_{\mathrm{HP}}$ open question]
\label{rem:lb_class_resolved}
Theorem~\ref{thm:spsc_lb_sharp} resolves the open question stated in
the original Remark~\ref{rem:lb_class_conf}: a non-projecting policy
\emph{cannot} beat $T^{2/3}$ in the regime $d\ge d_0 T^{1/3}$, because
the trace identity bounds its expected exploit-round KL contribution
by exactly the same $\|x_t\|^2/d\le R_\cA^2/d$ term that bounds a
projecting policy's. The intuition that ``off-subspace exploitation
along $v$ contributes additional KL'' is correct only for a
\emph{specific} adversarial $v$; averaging over $v$, the alignment
contributes $1/d$ per round regardless of whether the policy projects.
The LB therefore applies to the full class $\Pi$ (which contains all
algorithms in our experiments: LinUCB, D-LinUCB, SW-LinUCB, OFUL,
LinTS, LowOFUL, VOFUL, BOSS, Jedra, and SPSC).
\end{remark}

\begin{remark}[Scope and open directions]
\label{rem:open_multiseg}
Theorem~\ref{thm:spsc_lb_sharp} establishes the $T^{2/3}$
probe-acquisition lower bound for \emph{all} bounded-action policies,
matching SPSC's probe-term rate up to the $d^{2/3}$ prefactor gap
discussed in Remark~\ref{rem:dim_dep}. It does \emph{not} provide a
lower bound on the leading $r\sqrt{KT}$ exploit-regret term. Two
extensions are open:
\begin{enumerate}
\item \emph{Multi-segment lower bound for the probe-cost piece.} The
single-segment construction here corresponds to $K=1$. A $K$-segment
construction with adversarial $v_k$ per segment would extend the LB
to the piecewise-stationary setting; preliminary attempts indicate the
rate remains $\Omega(c^{1/3}T^{2/3})$ (i.e., $K$-independent), but the
adaptive-probe-count handling and full-trajectory TV-transfer require
careful additional argument.
\item \emph{Matching $\Omega(r\sqrt{KT})$ outside $\Pi_{\mathrm{HP}}$.}
Within $\Pi_{\mathrm{HP}}$ the standard linear-bandit minimax LB
\citep[Theorem~24.2]{lattimore2020bandit} can be applied to the
projected $r$-dim problem to argue that any policy that exploits in
the learned subspace incurs $\Omega(r\sqrt{KT})$ exploit regret;
making this fully formal requires accounting for the information gain
from isotropic Gaussian probes, which acts as a side-observation
channel.
\end{enumerate}
The empirical evidence of \S\ref{sec:experiments} indicates that
ambient-LinUCB-style methods not in $\Pi_{\mathrm{HP}}$ do incur
regret $\gtrsim r\sqrt{KT}$ across the $(d,r,K)$ grid, suggesting the
$r\sqrt{KT}$ rate is fundamental even outside $\Pi_{\mathrm{HP}}$.
A proof is beyond the scope of this work.
\end{remark}


% =====================================================================
% Bibliography note: the proofs cite two sources that may not yet
% be in references.bib.  Add the entries below if missing:
%
% @book{lattimore2020bandit,
%   title     = {Bandit Algorithms},
%   author    = {Lattimore, Tor and Szepesv{\'a}ri, Csaba},
%   publisher = {Cambridge University Press},
%   year      = {2020}
% }
%
% @book{tsybakov2009introduction,
%   title     = {Introduction to Nonparametric Estimation},
%   author    = {Tsybakov, Alexandre B.},
%   publisher = {Springer},
%   year      = {2009}
% }
% =====================================================================
